\SourceLeaf{7}
\begin{center}
THEORIE DER ANZIEHUNG\par
\bigskip
HOMOGENER ELLIPTISCHER SPHÄROIDISCHER KÖRPER\par
\bigskip
NACH EINER NEUEN METHODE BEHANDELT.\par
\medskip
\rule{.16\textwidth}{.3pt}
\end{center}
\Article{1}
\par
Es ist wohlbekannt, dass die genaue Bestimmung der Anziehung, die ein homogener elliptischer Sphäroidkörper auf einen beliebigen Punkt ausübt, zu den schwierigsten Aufgaben der physikalischen Astronomie gehört. Seit \textsc{Newton}s Zeiten haben sich mehrere Geometer wiederholt mit großem Eifer damit befasst. Zunächst war die Untersuchung auf ein Sphäroid beschränkt, das durch Drehung einer halben Ellipse um eine ihrer beiden Achsen entsteht. Der große \textsc{Newton} selbst zeigte, wie man die Anziehung findet, die ein Punkt auf der Achse erfährt. Zugleich gab er den Zusammenhang zwischen den Anziehungen an, die Punkte im Innern des Sphäroids auf demselben Durchmesser erfahren (\textit{Princip. Lib.} I. \textit{Prop.} XCI). Danach bestimmte der scharfsinnige \textsc{Mac Laurin} mit einer überaus eleganten Synthese die Anziehung auf Punkte, die auf der Oberfläche des Sphäroids oder in der Fortsetzung seiner Äquatorebene liegen. Damit war zugleich die Theorie der Anziehung auf Punkte im Innern des Sphäroids vollständig abgeschlossen, denn nach \textsc{Newton}s Lehrsatz ließ sich deren Fall leicht auf die Anziehung auf Punkte der Oberfläche zurückführen (De caussa physica fluxus et refluxus maris, in \textit{Recueil des pièces qui ont remporté les prix de l’acad. roi. des sc. T.} IV; \textit{Treatise of fluxions B.} I. \textit{Ch.} 14). Der bedeutende \textsc{Lagrange} zeigte später nicht minder elegant, wie sich das, was \textsc{Mac Laurin} durch Synthese herausgearbeitet hatte, durch Analysis gewinnen lässt. Zuvor waren ihr derartige Fragen unzugänglich erschienen. Damit war der Weg für weitere Fortschritte eröffnet (\textit{Nouv. Mém. de l’Acad. de Berlin} 1773). Es blieb nämlich noch die Anziehung auf Punkte außerhalb des Sphäroids zu bestimmen, die weder in der Verlängerung der Achse noch in der Fortsetzung des Äquators liegen,
